% Generated draft: requires mathematical and source review.
\documentclass[11pt]{article}
\usepackage[margin=25mm]{geometry}
\usepackage{fontspec}
\setmainfont{Latin Modern Roman}
\usepackage{amsmath,amssymb,mathtools,unicode-math}
\setmathfont{Latin Modern Math}
\usepackage{enumitem,needspace,xurl,hyperref}
\setlength{\parindent}{0pt}
\setlength{\parskip}{5pt}
\setlength{\emergencystretch}{3em}
\newcommand{\N}{\mathbb{N}}
\newcommand{\Z}{\mathbb{Z}}
\newcommand{\Q}{\mathbb{Q}}
\newcommand{\R}{\mathbb{R}}
\newcommand{\C}{\mathbb{C}}
\newcommand{\F}{\mathbb{F}}
\newcommand{\A}{\mathbb{A}}
\newcommand{\PP}{\mathbb{P}}
\newcommand{\E}{\mathbb{E}}
\newcommand{\eps}{\varepsilon}
\newcommand{\dd}{\,\mathrm{d}}
\begin{document}
\section*{UnsolvedMath ID: 4800005}
\textbf{Source problem code:} catalogue entry
\subsection{Definitions and mathematical statement}
In the usual existential recurrence convention, an integer sequence $r(n)$ is good for $q$-recurrence of powers of a single transformation if, for every invertible probability-preserving system $(X,\mathcal B,\mu,T)$ and every $A\in\mathcal B$ with $\mu(A)>0$, there is an index $n\geq1$ for which
\[
\mu\left(A\cap T^{-r(n)}A\cap T^{-2r(n)}A\cap\cdots\cap T^{-qr(n)}A\right)>0.
\]
In particular, good for $1$-recurrence means that $\mu(A\cap T^{-r(n)}A)>0$ for some $n$. ``Powers'' here refers to iterates of the same $T$, not independently chosen commuting transformations. The source does not reproduce its definition; this conventional interpretation requires checking against the cited survey.

Under that interpretation, the source asks whether, for every integer $\ell\geq1$ and every sequence $a(n)$ good for $\ell$-recurrence of powers, each sequence $a(n)^k$, for $k\in\{1,\ldots,\ell\}$, is good for $1$-recurrence. Explicitly, for every such $k$, system and positive-measure $A$, the requested conclusion is
\[
\exists n\geq1,\qquad \mu(A\cap T^{-a(n)^k}A)>0.
\]
The witnessing index may depend on $k$, the system, and $A$; no simultaneous witness for all $k$ is required.

The supplied summary reports a negative example for $\ell=2$, specifically for $k=2$, due to Griesmer (2024). It leaves the individual assertions for $\ell\geq3$ open. A negative $\ell=2$ case refutes the universal all-$\ell$ implication but does not by itself supply a sequence good for higher recurrence orders. Thus the solved registry label and the remaining parameter cases have different scope.
\subsection{Short English statement}
Does a sequence good for $\ell$-fold recurrence also give ordinary recurrence after raising its terms to each power $1,\ldots,\ell$? The supplied report gives a negative square example for $\ell=2$ and leaves higher recurrence orders unresolved.
\subsection{Sources}
\begin{itemize}
\item Local UnsolvedMath registry, permanent ID 4800005, exact source code catalogue entry, ``Multiple ergodic averages---Problem 5''; proposed by Nikos Frantzikinakis, 2016. Registry status solved, classification SOLVED-IN-LITERATURE.
\item Nikos Frantzikinakis, \emph{Some open problems on multiple ergodic averages} (2016), Problem 5, \url{https://arxiv.org/abs/1103.3808}. Supplied source-TeX extraction accessed July 29, 2026.
\item The supplied summary attributes the negative $\ell=2$, $k=2$ result to Griesmer (2024), without a title or precise publication reference. Provenance says the author's November 2025 progress page records a later result for this item but does not mark the full numbered prompt solved; no progress-page URL is supplied. Higher-$\ell$ status is preserved as a supplied assessment, not independently checked.
\item Provenance flags rights NEEDS\_REVIEW: public access was verified, but redistribution permission was not established.
\end{itemize}
\end{document}
